% Task 2 — Global, Continental, and National Interface Nodes
% Northeast Megaregion Freight Network Design
% Tonnage figures for global and national nodes are in thousand short tons (ktons).
% Continental throughput figures are BTS border-crossing volume proxies (vehicle crossings).

\section{Task 2 — Global, Continental, and National Interface Nodes}

The second task identifies the interface nodes through which freight enters and exits the
Northeast megaregion. These nodes bridge the regional hub network with broader global,
continental, and domestic logistics systems. Twenty-nine nodes are catalogued across three
tiers: nine global nodes (seaports and cargo airports), eight continental nodes
(US--Canada border crossings), and twelve national nodes (high-interaction external
domestic counties).

\subsection{External Flow Isolation and Bucket Assignment}

Interface nodes are anchored by external freight flows — those in which exactly one
endpoint lies within the 14-state NE study area. From the Task 1 O-D matrix, all
\textit{inbound} and \textit{outbound} records are retained; \textit{internal} and
\textit{outside} flows are excluded. The resulting external flow table is then classified
into four mutually exclusive buckets using a priority-ordered assignment rule:

\begin{enumerate}
  \item \textbf{Global maritime} — flows carried by water mode (FAF mode 3).
  \item \textbf{Continental border} — flows with one endpoint in a US--Canada border state
        (NY, VT, NH, or ME), after maritime flows are excluded.
  \item \textbf{Global air} — flows carried by air or combined truck-and-air modes (FAF
        modes 5 or 11), after the preceding two buckets are applied.
  \item \textbf{National domestic} — all remaining external flows not captured by the
        above rules.
\end{enumerate}

This bucket structure reflects the physical reality that freight entering or leaving the
Northeast does so through distinct infrastructure types: seaports handle bulk maritime
cargo; major international airports handle high-value air freight; the northern land border
mediates Canada--US overland exchange; and the national domestic tier covers the
megaregion's deep integration with Midwest, Mid-Atlantic, and Sun Belt supply chains.
Each bucket's aggregate FAF tonnage is then allocated proportionally across the shortlisted
nodes within that tier using infrastructure-specific throughput proxies.

\subsection{Global Interface Nodes: Seaports and Cargo Airports}

\paragraph{Maritime nodes.}
Major commercial seaports along the NE coastline are shortlisted from the NTAD Commercial
Strategic Seaports dataset. Two seaports are identified as primary global maritime
gateways: Hampton Roads (VA) and Philadelphia (PA). The global maritime bucket tonnage
is allocated equally between the two nodes, each receiving 20{,}205 ktons in 2025
(Table~\ref{tab:task2-global}). Hampton Roads is the dominant deep-water port complex on
the mid-Atlantic coast; Philadelphia anchies freight interoperability with the Delaware
River estuary.

\paragraph{Air cargo nodes.}
Seven major Northeast cargo airports are shortlisted from the NTAD Aviation Facilities
and NTAD Air-to-Truck Intermodal datasets. Throughput shares are derived from
each facility's estimated intermodal area (\texttt{EST\_AREA}), which serves as a
proxy for air cargo handling capacity. The global air bucket tonnage (479{,}243 ktons
in 2025) is then allocated across the seven airports in proportion to their area shares.

John F. Kennedy International Airport (JFK) dominates the air cargo tier, receiving
236{,}426 ktons (49.3\% of the air allocation), reflecting its role as the premier
international gateway for the New York metropolitan area. Philadelphia International
(PHL) follows at 92{,}050 ktons (19.2\%) and Newark Liberty International (EWR) at
73{,}353 ktons (15.3\%). Together, these three airports account for 84\% of NE air
cargo throughput.

\begin{table}[htbp]
\centering
\caption{Global interface nodes — allocated freight throughput, 2025 and 2030 (ktons).}
\begin{tabular}{llllrr}
\toprule
Node & State & Type & 2025 (ktons) & 2030 (ktons) \\
\midrule
\multicolumn{5}{l}{\textit{Seaports}} \\
Hampton Roads           & VA & Seaport       & 20{,}205 & 20{,}405 \\
Philadelphia            & PA & Seaport       & 20{,}205 & 20{,}405 \\
\midrule
\multicolumn{5}{l}{\textit{Cargo Airports}} \\
John F. Kennedy Intl    & NY & Cargo airport & 236{,}426 & 272{,}867 \\
Philadelphia Intl       & PA & Cargo airport &  92{,}050 & 106{,}237 \\
Newark Liberty Intl     & NJ & Cargo airport &  73{,}353 &  84{,}659 \\
Washington Dulles Intl  & DC & Cargo airport &  26{,}581 &  30{,}679 \\
Gen. E. L. Logan Intl   & MA & Cargo airport &  19{,}354 &  22{,}337 \\
Pittsburgh Intl         & PA & Cargo airport &  15{,}949 &  18{,}407 \\
BWI Thurgood Marshall   & MD & Cargo airport &  15{,}531 &  17{,}925 \\
\bottomrule
\end{tabular}
\label{tab:task2-global}
\end{table}

\subsection{Continental Interface Nodes: US--Canada Border Crossings}

US--Canada border crossings are sourced from the Bureau of Transportation Statistics
(BTS) Border Crossing Entry Data. The dataset is filtered to the US--Canada border and
to freight-relevant measurement categories (truck, truck container, rail container,
and rail crossings). Coverage is restricted geographically to New York, Vermont,
New Hampshire, and Maine — the four NE states with a direct Canadian land border.

All qualifying crossing records are aggregated to node level by summing total crossing
volume. The top eight nodes by crossing intensity are retained and ranked. The FAF
\textit{continental border} bucket tonnage is then allocated across these eight crossings
proportionally, using BTS crossing volume as the throughput proxy.

\begin{table}[htbp]
\centering
\caption{Continental interface nodes — BTS crossing volume proxy and relative share, 2025.}
\begin{tabular}{llrr}
\toprule
Border Crossing & State & Crossing Volume & Share (\%) \\
\midrule
Buffalo Niagara Falls   & NY & 30{,}162{,}118 & 49.2 \\
Champlain--Rouses Point & NY &  9{,}776{,}473 & 16.0 \\
Alexandria Bay          & NY &  6{,}417{,}140 & 10.5 \\
Derby Line              & VT &  3{,}204{,}800 &  5.2 \\
Highgate Springs        & VT &  3{,}179{,}169 &  5.2 \\
Houlton                 & ME &  2{,}991{,}741 &  4.9 \\
Jackman                 & ME &  2{,}777{,}799 &  4.5 \\
Calais                  & ME &  2{,}738{,}811 &  4.5 \\
\bottomrule
\end{tabular}
\label{tab:task2-continental}
\end{table}

As shown in Table~\ref{tab:task2-continental}, Buffalo Niagara Falls is the single
dominant Canada gateway, commanding 49.2\% of total NE--Canada crossing volume.
This reflects Buffalo's location at the western tip of New York State, directly across
from Ontario, on the most heavily used I-90 / QEW corridor. Champlain--Rouses Point
(16.0\%) and Alexandria Bay (10.5\%) rank second and third, capturing the Interstate
87 and Bridge of Islands corridors, respectively. The six eastern crossings in Vermont
and Maine together account for 24.3\% of total crossing volume, serving secondary
Quebec--Maritime exchange routes.

\subsection{National Interface Nodes: External Domestic Hubs}

The national tier identifies the external US counties that exchange the greatest
bilateral freight volumes with the Northeast. Both inbound external origins (counties
sending freight into the NE) and outbound external destinations (counties receiving NE
freight) are aggregated into a bilateral county interaction table using the
\textit{national domestic} bucket. Counties are ranked by 2025 aggregate bilateral
tonnage, and the top twelve are retained as national interface nodes. FAF
\textit{national domestic} bucket tonnage (213{,}438 ktons in 2025) is allocated across
the twelve nodes proportionally.

\begin{table}[htbp]
\centering
\caption{National interface nodes — allocated freight throughput, 2025 (ktons).}
\begin{tabular}{llr}
\toprule
County & State & 2025 (ktons) \\
\midrule
Lucas County (Toledo)          & OH & 24{,}970 \\
Harris County (Houston)        & TX & 24{,}363 \\
Cook County (Chicago)          & IL & 24{,}108 \\
Franklin County (Columbus)     & OH & 21{,}339 \\
Cuyahoga County (Cleveland)    & OH & 21{,}279 \\
Hamilton County (Cincinnati)   & OH & 21{,}234 \\
Los Angeles County             & CA & 15{,}785 \\
Webb County (Laredo)           & TX & 13{,}995 \\
Stark County (Canton)          & OH & 12{,}445 \\
Jefferson County (Louisville)  & KY & 11{,}356 \\
Mecklenburg County (Charlotte) & NC & 11{,}290 \\
Wayne County (Detroit)         & MI & 11{,}273 \\
\bottomrule
\end{tabular}
\label{tab:task2-national}
\end{table}

The national tier is strongly concentrated in Ohio, which contributes five of the twelve
nodes (Lucas, Franklin, Cuyahoga, Hamilton, and Stark counties). This reflects Ohio's
position as the primary logistics relay between the Northeast megaregion and the greater
Midwest freight system, with direct I-80, I-76, and I-70 corridor connections. Harris
County (Houston, TX) and Cook County (Chicago, IL) represent the two largest external
metros with deep supply-chain ties to the Northeast. Los Angeles County and Laredo, TX
capture Pacific Rim import flows and US--Mexico cross-border traffic, respectively.

\subsection{Summary of Interface Node Tier Structure}

Table~\ref{tab:task2-summary} consolidates the three interface tiers. In aggregate,
the 29 interface nodes represent the principal gateways through which the Northeast
megaregion exchanges freight with the rest of the world. These nodes are incorporated
directly into the composite demand and infrastructure map (Task 3.1) and inform the
external flow routing component of the network flow model (Task 8).

\begin{table}[htbp]
\centering
\caption{Summary of interface node tier structure.}
\begin{tabular}{llrr}
\toprule
Tier & Node Type & Node Count & Total Throughput 2025 \\
\midrule
Global    & Seaports + cargo airports  & 9  & 519{,}652 ktons \\
Continental & US--Canada border crossings & 8  & Proportional (BTS proxy) \\
National  & External domestic counties & 12 & 213{,}438 ktons \\
\midrule
\textbf{Total} & & \textbf{29} & \\
\bottomrule
\end{tabular}
\label{tab:task2-summary}
\end{table}

The global tier concentrates almost entirely along the I-95 spine and the major Northeast
airport hubs, with JFK and the Philadelphia--Hampton Roads seaport pair anchoring the
two maritime gateways. The continental tier is dominated by the Buffalo--Champlain axis.
The national tier's Ohio concentration highlights the I-80/I-76 Pennsylvania Turnpike
corridor as the primary inland freight artery connecting the Northeast to the broader
North American network. Together, these 29 nodes define the external boundary conditions
that any hyperconnected freight network design must respect.
